% Nouveau EGA V (cote 25): archive cote 25, site page 44 (PDF page 45).
% typescript. Draft compared: none.
% AI-transcribed by Claude Opus 5.5 (claude-opus-5-5, adaptive thinking, effort high), Message Batches API, batch ega25b (msgbatch_01A7n5YuXubUPNK1zsepi8RG), 2026-11-13; not yet checked by eye.
\scanpage{25/44}
\mspage{\hand{40}}

dans T tel que \hand{$\xi$} \var{\inserted{$\in U(k)$}}{$\in U(k)$}{insertion}{ega25_25_044-v0001} \var{\deleted{V}}{}{deletion}{ega25_25_044-v0002} implique qu'il existe exactement  un point \var{\xout{du}}{}{typewriter strike-out}{ega25_25_044-v0003} y de
$\underline{Y}^{\text{sing}}(k)$ au dessus de \hand{$\xi$}. Cela résultera de 8.\var{\over{\unsure{6}}{7}}{7}{writing over}{ega25_25_044-v0004}. c) qui implique plus
\var{\xout{généralement}}{}{typewriter strike-out}{ega25_25_044-v0005} précisément:

\uline{Corollaire} 8.\var{\over{\unsure{8}}{9}}{9}{writing over}{ega25_25_044-v0006}.  Supposons la condition (i) de 8.\var{\over{\unsure{7}}{8}}{8}{writing over}{ega25_25_044-v0007}. satisfaite, et
soit  \var{\xout{\unsure{T(k) un point tel que T soit lisse en} \illegible{1} U \unsure{Il en semble}}}{}{typewriter strike-out}{ega25_25_044-v0008}
U l'ouvert de T des points en lesquels T est lisse sur k. Alors \var{\inserted{$U\neq 0$, \var{\deleted{et}}{}{deletion}{ega25_25_044-v0010}}}{$U\neq 0$, \var{\deleted{et}}{}{deletion}{ega25_25_044-v0010}}{insertion}{ega25_25_044-v0009}
$\underline{Y}^{\text{sing}}|U \longrightarrow U$ est une immersion ouverte, \var{\xout{(et même un isomorphisme}}{}{typewriter strike-out}{ega25_25_044-v0011} \var{\inserted{injection $Y^{\text{sing}}|U$ ne contient pas des points de $\underline{Y}^{\unsure{\text{sing sing}}}$.}}{injection $Y^{\text{sing}}|U$ ne contient pas des points de $\underline{Y}^{\unsure{\text{sing sing}}}$.}{insertion}{ega25_25_044-v0012}
\var{\xout{si X est propre sur k donc $\underline{Y}^{\text{sing}}$xxxx T surjectif)}}{}{typewriter strike-out}{ega25_25_044-v0013} Si X est propre sur
k donc $g:\underline{Y}^{\text{sing}}\longrightarrow T$ surjectif, alors $\underline{Y}^{\text{sing}}|U\longrightarrow U$ est un isomorphi
me, et U est le plus grand ouvert de T ayant cette dernière propriété.

Tout d'abord, comme g est dominant et gén \var{\deleted{ramifié}}{}{deletion}{ega25_25_044-v0014} \var{\inserted{étale}}{étale}{insertion}{ega25_25_044-v0015}, il est
gén. étale, donc on peut \unsure{trouver|tmouver} un ouvert non vide V de T tel que
\var{\xout{\unsure{gén on peut supposer de plus g étale au(x)} \illegible{1} \unsure{que} \illegible{1} \unsure{ne contienne pas} $\underline{Y}$ \illegible{1}}}{}{typewriter strike-out}{ega25_25_044-v0016}
$\underline{Y}^{\text{sing}}|V \longrightarrow V$ soit étale et surjectif\var{\inserted{\var{\margin{ce qui implique V lisse$/k$}}{\placedmargin{ce qui implique V lisse$/k$}}{marginal note placed}{ega25_25_044-v0018}}}{\var{\margin{ce qui implique V lisse$/k$}}{\placedmargin{ce qui implique V lisse$/k$}}{marginal note placed}{ega25_25_044-v0018}}{insertion}{ega25_25_044-v0017}. Si alors \hand{$\xi$}$\in V(k)$ et si y est
un point de $\underline{Y}^{\text{sing}}(k)$ au dessus de \hand{$\xi$}, alors avec les notations de
8.\var{\over{\unsure{6}}{7}}{7}{writing over}{ega25_25_044-v0019}. c) \var{\xout{\unsure{Xlex}}}{}{typewriter strike-out}{ega25_25_044-v0020} l'espace H' n'est autre que l'espace tangent à T en \hand{$\xi$},
et comme on a noté à cet endroit, cela implique que le point x de X(k)
projection de y est déterminé comme le point orthogonal à H', donc est
uniquement déterminé, donc comme $\underline{Y}^{\text{sing}}\subset X\text{x}$\hand{$\check{P}$}, y est uniquement déter-%
miné. Cela prouve déjà que g est birationnel (étant \var{\inserted[typed]{dominant,}}{dominant,}{insertion}{ega25_25_044-v0021} \var{\xout{\unsure{non ramifié}}}{}{typewriter strike-out}{ega25_25_044-v0022}
génériquement \var{\deleted{non ramffié}}{}{deletion}{ega25_25_044-v0023} \var{\inserted{étale}}{étale}{insertion}{ega25_25_044-v0024} et gén radiciel). \var{\xout{\unsure{Pour Evidemment V H}}}{}{typewriter strike-out}{ega25_25_044-v0025}
D'autre part, le morphisme \hand{$\unsure{\nu|\gamma}$}(dont la définition en forme est évidente)
qui associe à tout \var{\over{x}{$\xi$}}{$\xi$}{writing over}{ega25_25_044-v0026}$\in U(k)$ l'unique point $x=$ \hand{$\nu$}(\hand{$\xi$}) $\in P$ orthogonal à
l'espace tangent à U en \hand{$\xi$}, coïncide sur V avec le composé
$V\rightarrow \underline{Y}^{\text{sing}}|V \longrightarrow X$, où la deuxième flèche est la projection; donc
\var{\inserted[typed]{posant}}{posant}{insertion}{ega25_25_044-v0027} \var{\xout{\unsure{La restriction de}}}{}{typewriter strike-out}{ega25_25_044-v0028} $h=($\hand{$\nu$}$,\mathrm{id}) : U\rightarrow P\text{x}T$ , \var{\xout{\unsure{(il induit une immersion à V et est}}}{}{typewriter strike-out}{ega25_25_044-v0029}
\var{\xout{\unsure{montre que le composé j g}}}{}{typewriter strike-out}{ega25_25_044-v0030} $g_1 = g|g^{-1}(U)$\var{\xout{=}}{}{typewriter strike-out}{ega25_25_044-v0031}: $g^{-1}(U)\rightarrow U$\var{\xout{\unsure{Ex}}}{}{typewriter strike-out}{ega25_25_044-v0032} , le composé

\var{\margin{44}}{\placedmargin{44}}{marginal note placed}{ega25_25_044-v0033}
